\begin{abstract}
Spectral optimizers like Muon achieve striking gains over AdamW in vision tasks, but what drives their advantage and when does it break down? We present a controlled empirical study comparing Muon across three architectures (ResNet-50, ConvNeXt-T, DeiT-S) and three datasets, with SOAP and Shampoo compared on CIFAR-100 (and Tiny-ImageNet in the appendix), that reveals that the spectral advantage holds across all tested architectures, datasets, and training lengths; for Muon, Newton--Schulz orthogonalization is the primary contributor. Removing it collapses performance by 17\,pp, erasing the entire spectral advantage and confirming that spectral preconditioning, not momentum or learning-rate scheduling, is the primary component. This orthogonalization reshapes weight spectra toward uniform singular-value distributions, and layer-wise spectral divergence between optimizers correlates with per-layer accuracy impact on CNNs, providing a mechanistic rather than purely empirical explanation. A normalization-swap experiment yields the most striking result: replacing batch normalization with layer normalization causes AdamW to collapse (67.85\% $\to$ 50.51\%) while Muon is unaffected (76.88\% $\to$ 78.52\%), demonstrating that spectral optimizers exhibit robustness to normalization changes compared to first-order methods. Muon converges 2--15$\times$ faster than AdamW in wall-clock time across CIFAR-100 and Tiny-ImageNet. These findings establish that spectral structure control is the primary driver of optimizer robustness across normalization regimes in modern vision architectures, particularly on convolutional architectures\footnote{Code is available at: \url{https://anonymous.4open.science/r/spectral_optim_vision_geometry-DC6A}}.
\end{abstract}
